%% toc: P15 164 164.3,164.4
\chapter{Popper on Probability and Quantum Theory}
\label{sec:p15-w02b-p15-027}

This chapter continues \cref{sec:p15-popper-hypothetico-deductive-reasoning}
with chapters~8 and~9 of \emph{The Logic of Scientific Discovery} and their
later notes, from the same Routledge Classics 2002 text (electronic issue
2005) of the 1959 translation. Its subject is the testing of probability
hypotheses: how a frequency theory of probability is built, why
probability statements cannot be falsified by finite records, which
methodological rule makes them testable, and how the same analysis is
applied to Heisenberg's uncertainty formulae. The inputs are frequency
hypotheses about event sequences; the conclusions are theorems of a
finite frequency calculus, a binomial formula, a Bernoulli bound, and
methodological and interpretive theses whose status is stated where they
occur. Corroboration of hypotheses, which Popper keeps separate from the
probability of events, is \cref{sec:p15-popper-corroboration}; the
rare-outcome theorem \cref{thm:p15-popper-coin} is the simplest case of
the decidability problem below.

\section{Probability statements and the finite calculus}
\label{sec:p15-w02b-p15-027-finite-frequency}

Chapter~8 of the primary text concerns the probability of events.
Whether a well-tested hypothesis is more probable than a little-tested
one is deferred to the corroboration sections. The chapter sets two
tasks: new foundations for the calculus, built as a frequency theory
after von Mises but without his limit axiom and with a weaker
randomness axiom; and an account of how probability statements relate
to experience, the problem of their decidability
\parencite[chapter~8, opening]{Popper2002LogicDiscoveryUI}.

Section~47 separates numerical probability statements, such as a value
of 1/18 for throwing eleven with two true dice, from non-numerical ones.
Section~48 finds three readings latent in the classical ratio of
favourable to equally possible cases. A psychologistic reading measures
feelings of certainty. A logical variant, chiefly Keynes's, treats
probability as a relation between statements whose extremes are
derivability (value 1) and contradiction (value 0). The objective
reading takes a numerical probability statement to assert a relative
frequency of a kind of event within a sequence of occurrences. Popper
adopts a modified frequency theory because, on his argument, only an
objective theory explains the use of the calculus in empirical science.
The subjective theory decides probability statements consistently, but
only by making them tautologies. He also rejects deriving objective
frequency statements from subjective premises with Bernoulli's theorem
as a bridge \parencite[\S\S47--48]{Popper2002LogicDiscoveryUI}.
His later note to that passage keeps the rejection while allowing the
theorem to bridge propensities and statistics
\parencite[\S48, note~6]{Popper2002LogicDiscoveryUI}.

Section~49 names the \emph{fundamental problem of the theory of chance}:
apparently incalculable events are just the ones to which the calculus
is applied, and the frequency conclusions drawn from that ignorance are
corroborated in practice. A subjective theory removes the air of paradox
only by making the calculus a transformation of what we do not know,
which does not explain empirical corroboration
\parencite[\S49]{Popper2002LogicDiscoveryUI}.

Von Mises defines a \emph{collective} as an indefinitely continuable
event sequence satisfying an axiom of convergence (the relative
frequency of each property tends to a limit) and an axiom of randomness
(no gambling system changes those limits in a selected subsequence).
Probability is the limit of relative frequency in a collective, and the
calculus derives new distributions from given ones
\parencite[\S50]{Popper2002LogicDiscoveryUI}.
Section~51 holds the convergence axiom no less responsible for the
difficulties than the randomness axiom. One task is mathematical: improve
the randomness axiom so that it yields the binomial formula in a third
form, from which Bernoulli's theorem follows. The other is epistemological:
eliminate the convergence axiom. Three symbols mark the stages: \(F''\)
for relative frequency in finite classes, \(F'\) for the limit of an
infinite frequency sequence, and \(F\) for objective probability, the
relative frequency in a chance-like sequence
\parencite[\S51]{Popper2002LogicDiscoveryUI}.

\begin{hbdefinition}{def:p15-w02b-p15-027-finite-frequency}{Relative frequency in a finite reference class}
Let \(\alpha\) be a nonempty finite class, the \emph{reference class},
and let \(\beta,\gamma\) be arbitrary classes, the \emph{property
classes}. Write \(N(\alpha)\) for the number of elements of \(\alpha\),
\(\alpha.\beta\) for the product class of \(\alpha\) and \(\beta\),
\(\beta+\gamma\) for the class of elements belonging to \(\beta\) or to
\(\gamma\) (non-exclusive), and \(\bar\beta\) for the complement of
\(\beta\). Popper's Definition~1 is
\[
  {}_{\alpha}F''(\beta)=\frac{N(\alpha.\beta)}{N(\alpha)} .
\]
When \(N(\alpha.\beta)>0\), the class \(\alpha.\beta\) is itself a
reference class, the \emph{selection from \(\alpha\) according to
\(\beta\)}. Properties \(\beta\) and \(\gamma\) are \emph{mutually
independent within \(\alpha\)} when
\({}_{\alpha.\beta}F''(\gamma)={}_{\alpha}F''(\gamma)\); equivalently
\(\gamma\) is \emph{insensitive} within \(\alpha\) to selection of
\(\beta\)-elements.
No assumption that the elements of \(\alpha\) are equally possible is
made \parencite[\S\S52--53, with note~*1 to \S52]{Popper2002LogicDiscoveryUI}.
The positivity condition on \(N(\alpha.\beta)\) is stated here because a
selected class must be a nonempty reference class.
\end{hbdefinition}

Popper stresses that the theorems derived from Definition~1 contain only
relative frequencies; class cardinalities appear only in the proofs
of the first theorems \parencite[\S52]{Popper2002LogicDiscoveryUI}.
Appendix~ii lists the finite calculus.

\begin{hbtheorem}{thm:p15-w02b-p15-027-finite-calculus}{proved=annex:p15-w02b-p15-027-finite-calculus}{Popper's finite frequency calculus (appendix~ii)}
Let \(\alpha\) be a nonempty finite class and \(\beta,\gamma\) classes.
\begin{enumerate}
\item \({}_{\alpha}F''(\beta)+{}_{\alpha}F''(\bar\beta)=1\).
\item General multiplication theorem: if \(N(\alpha.\beta)>0\),
  respectively \(N(\alpha.\gamma)>0\), then
  \[
   {}_{\alpha}F''(\beta.\gamma)={}_{\alpha}F''(\beta)\cdot{}_{\alpha.\beta}F''(\gamma),
  \]
  respectively
  \[
   {}_{\alpha}F''(\beta.\gamma)={}_{\alpha.\gamma}F''(\beta)\cdot{}_{\alpha}F''(\gamma).
  \]
\item If \(N(\alpha.\beta)>0\) and \(N(\alpha.\gamma)>0\), the following
  are equivalent: \(\gamma\) is independent of \(\beta\) within \(\alpha\);
  \(\beta\) is independent of \(\gamma\) within \(\alpha\); the special
  multiplication theorem
  \({}_{\alpha}F''(\beta.\gamma)={}_{\alpha}F''(\beta)\cdot{}_{\alpha}F''(\gamma)\).
\item General addition theorem:
  \({}_{\alpha}F''(\beta+\gamma)={}_{\alpha}F''(\beta)+{}_{\alpha}F''(\gamma)
  -{}_{\alpha}F''(\beta.\gamma)\); if \(N(\alpha.\beta.\gamma)=0\) the last
  term vanishes.
\item General division theorem: if \(N(\alpha.\beta)>0\) then
  \[
   {}_{\alpha.\beta}F''(\gamma)={}_{\alpha}F''(\beta.\gamma)/{}_{\alpha}F''(\beta).
  \]
\item Fourth form of Bayes's theorem: let \(\beta_1,\ldots,\beta_n\)
  satisfy \(N(\alpha.\beta_i.\beta_j)=0\) for \(i\neq j\) and
  \(N(\alpha.\beta_i)>0\) for every \(i\); let
  \(\alpha.\gamma\subset\Sigma\beta_i\) and \(N(\alpha.\gamma)>0\). Then
  for every \(i\)
  \[
   {}_{\alpha.\gamma}F''(\beta_i)
   =\frac{{}_{\alpha.\beta_i}F''(\gamma)\cdot{}_{\alpha}F''(\beta_i)}
          {\sum_{j=1}^n{}_{\alpha.\beta_j}F''(\gamma)\cdot{}_{\alpha}F''(\beta_j)} .
  \]
\item Let \(\gamma\subset\beta\). Then
  \({}_{\alpha}F''(\gamma)\leq{}_{\alpha}F''(\beta)\), and equality
  holds exactly when \(N(\alpha.\beta.\bar\gamma)=0\).
\end{enumerate}
\end{hbtheorem}

Items 1--6 are Popper's formulas (1), (1\('\)), (1\(_s\)), (2), (2\(_s\)),
(3) and (4\(_s\)), with \(\Sigma\) placed before class names for their
union \parencite[\S52; appendix~ii]{Popper2002LogicDiscoveryUI}.
The positivity hypotheses are added here so that every selected
reference class is nonempty; appendix~ii leaves them tacit. Item~7 is the
finite counterpart of the subclass comparison used in \S72 and is a
handbook addition. Popper reads item~3 as showing that independence is
symmetric in the two property classes, and the division theorem as
showing that independence is a special case of selection.

A 1934 footnote claimed more: that independence is symmetric in all
three classes provided \(\beta\) and \(\gamma\) are finite. The later
starred addition to that footnote declares the finiteness condition
insufficient and gives a counterexample. It suggests the intended
condition was that \(\alpha\) bound \(\beta\) and \(\gamma\), or that
\(\alpha\) be the universe of discourse
\parencite[\S53, note~4]{Popper2002LogicDiscoveryUI}.

\begin{hbtheorem}{thm:p15-w02b-p15-027-triple-symmetry}{proved=annex:p15-w02b-p15-027-sym}{Popper's five-button counterexample}
Let a universe of five buttons contain four round ones (\(\alpha\)),
two round and black (\(\alpha\beta\)), two round and large
(\(\alpha\gamma\)), one round, black and large (\(\alpha\beta\gamma\)),
and one square, black and large (\(\bar\alpha\beta\gamma\)).
Then \(\beta\) and \(\gamma\) are mutually independent within
\(\alpha\), and
\({}_{\beta.\alpha}F''(\gamma)={}_{\alpha}F''(\gamma)=1/2\neq2/3={}_{\beta}F''(\gamma)\).
So \(\alpha\) and \(\gamma\) are not independent within \(\beta\).
If instead \(\beta\subset\alpha\), then
\({}_{\beta.\alpha}F''(\gamma)={}_{\beta}F''(\gamma)\) for every \(\gamma\).
\end{hbtheorem}

The source states the counterexample as
\({}_{\alpha}F''(\gamma)\neq{}_{\beta}F''(\gamma)\). Reading
\emph{triple symmetry} as the further claim that independence of
\(\beta,\gamma\) within \(\alpha\) yields independence of
\(\alpha,\gamma\) within \(\beta\) is this handbook's interpretation;
under it, the displayed inequality is exactly the failure.

\section{After-effect, \texorpdfstring{\(n\)}{n}-freedom and the first binomial formula}
\label{sec:p15-w02b-p15-027-n-freedom}

Section~53 reads independence in subjective vocabulary as irrelevance:
information that an element has \(\beta\) is irrelevant to whether it
has \(\gamma\) exactly when the two are independent. A later note adds
that relevance can be defined from dependence, but not conversely
\parencite[\S53, note~*1]{Popper2002LogicDiscoveryUI}.
Section~54 numbers a finite reference class and distinguishes two
selections: \emph{ordinal selection}, by a property depending on an
element's ordinal number, and \emph{neighbourhood selection}, by
properties of its predecessors or successors. \emph{Primary properties}
such as heads belong to an element regardless of position; \emph{secondary
properties} such as being even-numbered or being the successor of tails
come from position. A sequence with two primary properties, written
\(1\) and \(0\), is an \emph{alternative}
\parencite[\S54]{Popper2002LogicDiscoveryUI}.

Section~55 calls an alternative \emph{1-free} when its primary-property
frequencies are insensitive to selection according to a single
predecessor, \emph{2-free} when also insensitive to pairs of
predecessors, and \emph{\(n\)-free} when insensitive to selection
according to every \(k\)-tuple of predecessors for \(k\leq n\).
Popper uses the generating periods
\[
\begin{aligned}
(A)&\ 1100, &\qquad (B)&\ 10111000,\\
(C)&\ 1011000011110100, &\qquad
(D)&\ 01100011101010010000010111110011
\end{aligned}
\]
for 1-, 2-, 3- and 4-free alternatives with equal distribution. He
states that such a period has at least \(2^{n+1}\) elements, and that
the period followed by its first \(n\) elements contains every
\((n+1)\)-tuple of zeros and ones. He neglects deviations caused by the
first and last terms of a finite sequence
\parencite[\S55]{Popper2002LogicDiscoveryUI}.
A footnote credits K.~Schiff with the remark that insensitivity to
predecessor \(n\)-tuples already implies insensitivity to shorter ones
\parencite[\S55, note~1]{Popper2002LogicDiscoveryUI}.
The term \emph{free from after-effect} is Reichenbach's, while the
recursive graded concept is claimed as Popper's; a note added in 1968
reports an earlier use of the term by Smoluchowski
\parencite[\S55, note~*2]{Popper2002LogicDiscoveryUI}.
Reichenbach's own treatment of after-effect and normal sequences is
in \cref{sec:p15-reichenbach-normality}.

To state these claims exactly, the handbook reads a finite period
cyclically, as one turn of an endlessly repeated period. Boundary
deviations then vanish instead of being neglected.

\begin{hbdefinition}{def:p15-w02b-p15-027-cyclic-freedom}{Cyclic counts and \texorpdfstring{\(n\)}{n}-freedom of a finite period}
Let \(N\geq1\) and let \(a=(a_0,\ldots,a_{N-1})\) be a finite alternative,
indices read modulo \(N\). For a word \(w=w_1\cdots w_k\) over
\(\{0,1\}\) with \(k\geq1\), let \(c_a(w)\) be the number of
\(i\in\{0,\ldots,N-1\}\) such that \(a_{i+j-1}=w_j\) for
\(j=1,\ldots,k\). Put \(p=c_a(1)/N\) and \(q=c_a(0)/N\).
The positions \(t\) whose \(k\) immediate predecessors form the word
\(\tau\) are \(c_a(\tau)\) in number, and \(c_a(\tau1)\) of them carry
the property \(1\).
Call \(a\) \emph{cyclically \(n\)-free} (\(n\geq0\)) if
\(c_a(\tau1)=p\,c_a(\tau)\) for every word \(\tau\) of length
\(k\in\{1,\ldots,n\}\) with \(c_a(\tau)>0\).
The \emph{sequence of overlapping \(n\)-segments} \(a_{(n)}\) consists
of the \(N\) cyclic windows of length \(n\); for \(m\leq n\),
\({}_{a_{(n)}}F''(m)\) is the proportion of those windows containing
exactly \(m\) ones, and \({}_{a_{(n)}}F''(\sigma_m)=c_a(\sigma_m)/N\)
for one fixed arrangement \(\sigma_m\) of length \(n\) with \(m\) ones.
\end{hbdefinition}

The counts \(c_a\) and the cyclic reading are handbook devices; the
symbols \(\alpha_{(n)}\), \(\sigma_m\) and \({}_{\alpha_{(n)}}F''(m)\) are
Popper's \parencite[\S56; appendix~iii]{Popper2002LogicDiscoveryUI}.

\begin{hbtheorem}{thm:p15-w02b-p15-027-binomial}{proved=annex:p15-w02b-p15-027-binomial}{Schiff's reduction and the first binomial formula, cyclic version}
Let \(a\) be a finite alternative of length \(N\), with \(p,q\) as in
\cref{def:p15-w02b-p15-027-cyclic-freedom}.
\begin{enumerate}
\item If \(n\geq1\) and \(c_a(\rho1)=p\,c_a(\rho)\) for every word
  \(\rho\) of length exactly \(n\) with \(c_a(\rho)>0\), then \(a\) is
  cyclically \(n\)-free.
\item If \(n\geq1\) and \(a\) is cyclically \((n-1)\)-free, then every
  word \(w\) of length \(k\leq n\) with \(m\) ones satisfies
  \(c_a(w)=N\,p^m q^{k-m}\). Hence
  \({}_{a_{(n)}}F''(\sigma_m)=p^mq^{n-m}\) and
  \[
    {}_{a_{(n)}}F''(m)=\binom{n}{m}p^mq^{n-m}\qquad(0\leq m\leq n).
  \]
\item If \(a\) is cyclically \(n\)-free and \(p=q=1/2\), then every word
  of length \(n+1\) occurs exactly \(N/2^{n+1}\) times. Hence
  \(2^{n+1}\) divides \(N\), and the linear word
  \(a_0\cdots a_{N-1}a_0\cdots a_{n-1}\) of length \(N+n\) contains every
  word of length \(n+1\) as a block of consecutive terms.
\item Conversely, if \(N=2^{n+1}\) and every word of length \(n+1\)
  occurs exactly once, then \(p=q=1/2\) and \(a\) is cyclically \(n\)-free.
\end{enumerate}
\end{hbtheorem}

Item~2 is appendix~iii's induction with the boundary error removed.
Popper's linear statement assumes \(\alpha\) at least \((n-1)\)-free
and neglects errors at the last term
\parencite[appendix~iii]{Popper2002LogicDiscoveryUI}. Item~3 gives his
lower bound \(2^{n+1}\) and his \((n+1)\)-tuple property for equal
distribution. Item~4 is a handbook converse used to check his periods.

\begin{hbexample}{def:p15-w02b-p15-027-periods}{Popper's four generating periods}
Read cyclically, period (A) has the windows \(11,10,00,01\) of length
2, each once; by item~4 it is cyclically 1-free with \(p=1/2\). It is
not 2-free: the only element whose two predecessors are \(1,1\) is a
zero, so selection by the pair \(11\) gives frequency \(0\neq1/2\) for
ones. Popper's own observation is that all four pair-selections give
uninterrupted blocks \parencite[\S55]{Popper2002LogicDiscoveryUI}.
Period (B) has the eight windows
\(101,011,111,110,100,000,001,010\) of length 3, each once, so it is
cyclically 2-free. Periods (C) and (D) contain each word of length 4,
respectively 5, exactly once (listed in the annex), so they are
cyclically 3-free and 4-free. By item~2 with \(n=3\), in (B) each of the
eight arrangements of three terms has frequency \(1/8\) and
\({}_{a_{(3)}}F''(m)=\binom{3}{m}/8\).
\end{hbexample}

The periods are generated by an explicit rule and are nevertheless
\(n\)-free. This is the finite form of Popper's later point that
chance-like sequences need not lack a mathematical law (\S69).

\section{Randomness and Bernoulli's theorem}
\label{sec:p15-w02b-p15-027-bernoulli}

Section~56 introduces \(n\)-segments of a sequence \(\alpha\): the
sequence \(\alpha_{(n)}\) of overlapping segments (elements \(x\) to
\(x+n-1\) for each \(x\)) and a sequence \(\alpha_n\) of adjoining
segments (non-overlapping blocks). For a finite alternative at least
\((n-1)\)-free with \({}_{\alpha}F''(1)=p\) and \({}_{\alpha}F''(0)=q\), the
first binomial formula
\({}_{\alpha_{(n)}}F''(m)={}^nC_m\,p^mq^{n-m}\) holds; its proof is
appendix~iii (\cref{thm:p15-w02b-p15-027-binomial}). The corresponding
problem for infinite sequences of adjoining segments is
\emph{Bernoulli's problem}, and for overlapping ones the
\emph{quasi-Bernoulli problem}
\parencite[\S56 and note~1]{Popper2002LogicDiscoveryUI}.

Section~57 passes to infinite reference sequences, provisionally using
the limit \(F'\) of relative frequencies. A \emph{mathematical sequence}
is fixed by a mathematical rule, such as zero exactly at positions
divisible by four, giving \({}_{\alpha}F'(1)=3/4\); an \emph{empirical
sequence} is fixed by an empirical rule, such as the outcomes of tosses
of a given coin. The distinction is intensional: no finite segment
settles which kind of sequence it belongs to. Restricting attention to
mathematical sequences with convergent frequency sequences introduces
an axiom of convergence. To conjecture that a mathematical sequence
approximates an empirical one in its frequencies is to frame a
frequency hypothesis \parencite[\S57]{Popper2002LogicDiscoveryUI}.
Such estimates have two typical sources: equal-chance hypotheses, often
from symmetry, and extrapolation of counted statistics, as with
mortality rates. Neither is a logical derivation from counted cases;
every predictive estimate about an infinite empirical sequence goes
beyond the observations. Popper avoids the labels \emph{a priori} and
\emph{a posteriori} for the two sources
\parencite[\S57]{Popper2002LogicDiscoveryUI}.

Section~58 accepts Kamke's objection that a collective in von Mises's
sense cannot be exhibited, because any rule defining an example is a
gambling system. Excluding every gambling system is also more than
Bernoulli's theorem needs. Popper replaces the randomness axiom by
\emph{absolute freedom}: \(n\)-freedom for every \(n\). This excludes
selection by a constant neighbourhood characteristic, selection of
elements at a constant distance apart, and combinations of these:
systems usable without knowing the first element. Selections tied to
absolute position, such as prime-numbered places, are not excluded.
Because absolutely free sequences can be constructed by rule
(appendix~iv), the concept is not empty
\parencite[\S58]{Popper2002LogicDiscoveryUI}.
Later notes record that the last seven words of his reformulation of von
Mises's axiom are absent from the German text, and that excluding a given
denumerable set of gambling systems is compatible with constructing an
example, by Wald's diagonal method
\parencite[\S58, notes~*1, *2]{Popper2002LogicDiscoveryUI}.

\begin{hbdefinition}{def:p15-w02b-p15-027-absolute-freedom}{Chance-like sequences (\S59)}
Let \(a=(a_1,a_2,\ldots)\) be an infinite alternative whose ones have the
frequency limit \({}_{a}F'(1)=p\). For a word \(\tau\) of length
\(k\geq1\), the predecessor selection \(S_\tau\) is the set of positions
\(t>k\) with \((a_{t-k},\ldots,a_{t-1})=\tau\). The sequence is
\emph{absolutely free} (\emph{chance-like}, \emph{random}) if for every
\(k\) and every \(\tau\) for which \(S_\tau\) is infinite, the relative
frequency of ones along \(S_\tau\) converges to \(p\). Its objective
probability is then \({}_{a}F(1)={}_{a}F'(1)=p\)
\parencite[\S59]{Popper2002LogicDiscoveryUI}.
\end{hbdefinition}

The restriction to infinite selections is this handbook's; a finite
selection has no frequency limit. A later note says Popper would now use
\emph{objective probability} more widely, for every objective
interpretation of the formal calculus, the propensity interpretation
included; in \S59 it is an auxiliary concept of one frequency theory
\parencite[\S59, note~*1]{Popper2002LogicDiscoveryUI}.

Section~60 states three binomial formulas:
\[
\begin{aligned}
 {}_{\alpha_{(n)}}F''(m)&={}^nC_m\,p^mq^{n-m} &&\text{(1) finite, overlapping},\\
 {}_{\alpha_{(n)}}F'(m)&={}^nC_m\,p^mq^{n-m} &&\text{(2) infinite, overlapping},\\
 {}_{\alpha_{n}}F(m)&={}^nC_m\,p^mq^{n-m} &&\text{(3) infinite, adjoining}.
\end{aligned}
\]
Formula (2) needs only that \(\alpha\) be at least \((n-1)\)-free. Formula
(3), the solution of Bernoulli's problem, also asserts that sequences of
adjoining segments are themselves chance-like. Popper gives a two-step
argument, which he calls a sketch. First, adjoining segments are a
normal ordinal selection from overlapping ones. If a pair such as
\(0,0\) were too frequent in one adjoining sequence, it would recur
preferentially at characteristic distances in \(\alpha\), contrary to
absolute freedom. Second, the adjoining-segment sequences are
absolutely free by a parallel compensation argument
\parencite[\S60 and note~*1]{Popper2002LogicDiscoveryUI}.
This handbook does not prove formula (3) from
\cref{def:p15-w02b-p15-027-absolute-freedom}; the two steps remain
Popper's sketch. A later note withdraws the claim that absolute freedom
gives insensitivity to \emph{all} combinations of ordinal and
neighbourhood selection. It restricts the claim to those usable as
gambling systems, after a correction by Wald in 1935
\parencite[\S60, note~*3]{Popper2002LogicDiscoveryUI}.
Note~2 to \S60 records that Reichenbach denied that freedom from
after-effect yields his normal sequences, those satisfying formula (3);
Popper's argument asserts that it does.

Section~61 calls an \(n\)-segment \(\Delta p\) when its frequency of
ones \(m/n\) satisfies \(|m/n-p|<\delta\), for a fixed \(\delta>0\).
Bernoulli's theorem is
\(\lim_{n\to\infty}{}_{\alpha_n}F(\Delta p)=1\). It follows from
formula~(3) by arithmetic, and the same method applied to (2) gives
\(\lim_{n\to\infty}{}_{\alpha_{(n)}}F'(\Delta p)=1\), a variant for
overlapping segments and sequences with after-effect. In frequency
language: almost all sufficiently long segments are \emph{fair samples}
whose frequency lies close to \(p\). Short segments fluctuate; long ones
behave quasi-convergently, which Popper calls statistical stability
\parencite[\S61]{Popper2002LogicDiscoveryUI}.

\begin{hbtheorem}{thm:p15-w02b-p15-027-bernoulli-bound}{proved=annex:p15-w02b-p15-027-bernoulli-bound}{The arithmetic behind Bernoulli's theorem}
Let \(0\leq p\leq1\), \(q=1-p\), \(n\geq1\) and \(\delta>0\). Then
\[
 \sum_{\substack{0\leq m\leq n\\ |m/n-p|\geq\delta}}\binom{n}{m}p^mq^{n-m}
 \;\leq\;\frac{pq}{n\delta^2}.
\]
Consequently:
\begin{enumerate}
\item If a finite alternative \(a\) is cyclically \((n-1)\)-free, the
  proportion of its overlapping cyclic \(n\)-segments that are not
  \(\Delta p\) is at most \(pq/(n\delta^2)\).
\item If formula (3) holds for every \(n\), then
  \({}_{\alpha_n}F(\Delta p)\geq1-pq/(n\delta^2)\), so
  \(\lim_{n\to\infty}{}_{\alpha_n}F(\Delta p)=1\).
\item Let \(0<p<1\). Given \(\varphi>0\) and \(0<\varepsilon_0\leq1\), if
  \(n\geq pq/(\varepsilon_0\varphi^2)\) then for every
  \(\varepsilon\in[\varepsilon_0,1]\) the bound
  \(\delta_\varepsilon=\sqrt{pq/(n\varepsilon)}\), at which the
  displayed sum is at most \(\varepsilon\), satisfies
  \(\delta_\varepsilon\leq\varphi\).
\end{enumerate}
\end{hbtheorem}

The inequality is Chebyshev's bound for binomial weights. Popper cites
no particular proof and refers to textbooks
\parencite[\S61]{Popper2002LogicDiscoveryUI}. Item~1 makes exact his
\S63 remark that the theorem holds approximately for finite sequences
that are \(n\)-free for large \(n\), without any convergence axiom.
Item~3 is a handbook illustration for \S68, not his analysis of a
characteristic value of \(\Delta p\). A general weak law in Feller's form
is \cref{sec:p15-feller-generalized-weak-law}.

Section~62 argues that the subjective theory cannot read Bernoulli's
theorem as the law of great numbers: \(p\) would mean a degree of
rational belief before the deduction and a frequency after it. Only
frequencies can deviate from frequencies
\parencite[\S62 and note~*1]{Popper2002LogicDiscoveryUI}.
Section~63 cites von Mises's observation that some sequences have a
frequency limit and still contain, with frequency near 1, long segments
deviating greatly from it. The final digits of a six-figure table of
square roots are his example. So convergence does not yield Bernoulli's
theorem, and the derivation rests on absolute freedom; Popper conjectures
that the theorem is independent of the convergence axiom
\parencite[\S63 and note~1]{Popper2002LogicDiscoveryUI}. A later note
keeps the doubt about the convergence axiom and refers to appendix~*vi
and to Borel's normal numbers. It no longer endorses the view implicit in the sentence that follows
it, which proposes to show the theorem also for mathematical sequences
whose primary properties have no frequency limits
\parencite[\S63, note~*1]{Popper2002LogicDiscoveryUI}.

Section~64 replaces the frequency limit by a \emph{middle frequency}: an
accumulation point of the sequence of relative frequencies
\parencite[\S64]{Popper2002LogicDiscoveryUI}.

\begin{hbtheorem}{thm:p15-w02b-p15-027-middle-frequencies}{proved=annex:p15-w02b-p15-027-middle-frequencies}{Middle frequencies always exist and form an interval}
Let \(a=(a_1,a_2,\ldots)\) be an infinite alternative and
\(h_n=(a_1+\cdots+a_n)/n\). Then:
\begin{enumerate}
\item the set \(M\) of accumulation points of \((h_n)\), the middle
  frequencies, is nonempty and lies in \([0,1]\);
\item \(M=[\liminf_n h_n,\limsup_n h_n]\), so if \(M\) has more than one
  element it is a nondegenerate interval;
\item \(M\) has exactly one element if and only if \((h_n)\) converges,
  and then that element is the frequency limit.
\end{enumerate}
\end{hbtheorem}

Item~1 is Popper's appeal to the Bolzano--Weierstrass theorem, item~3
his remark that a unique middle frequency is the frequency limit and
that a sequence without a limit has more than one, and item~2 his
note~2, which says without proof that several middle frequencies form a
continuum \parencite[\S64 and note~2]{Popper2002LogicDiscoveryUI}.

Middle frequencies need not be unique, so a definition of objective
probability through them alone would be ambiguous. Popper therefore
defines: an alternative is chance-like when the ones have exactly one
absolutely free middle frequency \(p\), its objective probability. He
splits this into a \emph{requirement of randomness} (at least one
absolutely free middle frequency) and a \emph{requirement of uniqueness}
(only one probability for one property in one alternative). The old
system adds the axiom of convergence (2\('\)): no further middle
frequency apart from \(p\). He refers to appendix~iv for a sequence with
one probability and no frequency limit, and claims Bernoulli's theorem
and the classical calculus follow from the two requirements
\parencite[\S64]{Popper2002LogicDiscoveryUI}. The fundamental problem
is then answered: the inference from irregularity to stability is valid
once irregularity is expressed as the hypothesis that one middle
frequency is free from after-effect. The older frequency theory had
postulated stability in the large through the convergence axiom
\parencite[\S64]{Popper2002LogicDiscoveryUI}. A later note offers a
\emph{requirement of finitude} (+2) in place of uniqueness: the sequence
should become \(n\)-free as quickly as possible and for the largest
possible \(n\), being approximately a shortest random-like sequence
\parencite[\S64, note~*2]{Popper2002LogicDiscoveryUI}. Popper's summary
note says his differences from von Mises concern only the limit axiom,
replaced by uniqueness, and the randomness axiom, modified so that model
sequences can be built \parencite[\S64, note~7]{Popper2002LogicDiscoveryUI}.

\section{Decidability and the logical form of probability statements}
\label{sec:p15-w02b-p15-027-decidability}

Section~65 states that while the binomial formula is derivable,
probability statements are not falsifiable: no basic statement or finite
conjunction of them contradicts an estimate. A long greatly deviating
segment has small but positive probability and is even to be expected;
a frequent occurrence of such segments can be redescribed as one
occurrence of a longer deviating segment. Only an intensionally defined
infinite sequence could contradict the estimate. Probability hypotheses
thus have infinite dimension and would be empirically uninformative,
though not void of logical content. Physics nevertheless uses them, and
a physicist judges when one is practically falsified. That practice
needs a methodological decision to treat highly improbable events as
prohibited \parencite[\S65 and note~1]{Popper2002LogicDiscoveryUI}.

\begin{hbtheorem}{thm:p15-w02b-p15-027-finite-records}{proved=annex:p15-w02b-p15-027-finite-records}{Finite records against chance and all-and-some hypotheses}
\begin{enumerate}
\item Let \(0\leq p\leq1\). If some infinite alternative is absolutely free with
  probability \(p\) (\cref{def:p15-w02b-p15-027-absolute-freedom}), then
  for every finite word \(w\) there is an absolutely free alternative
  with probability \(p\) whose first \(|w|\) terms are \(w\).
\item On infinite alternatives \(a=(a_1,a_2,\ldots)\) let
  \(S_\infty\) say that for every \(x\) there is \(y>x\) with \(a_y=0\),
  and \(S_b\) that for every \(x\), \(a_x=1\) implies \(a_{x+1}=0\).
  Every finite word is an initial segment both of an alternative
  satisfying \(S_\infty\) and of one violating it. Every alternative
  beginning \(1,1\) violates \(S_b\).
\end{enumerate}
\end{hbtheorem}

Item~1 is conditional on existence, which Popper supports through the
constructions of appendix~iv; those constructions are not reproduced
here. It shows that no initial record contradicts the hypothesis that
the sequence is chance-like with probability \(p\). Item~2 formalizes a
later note to \S66. What makes a statement of the form \emph{for every
\(x\) there is a \(y\)} unfalsifiable is that the \(y\) may be delayed
without bound; with the distance bounded, as in \(S_b\) (Popper's own
penny example), the statement can be falsifiable and indeed falsified.
He adds that his unguarded 1934 wording had been misread as a general
criticism of falsifiability \parencite[\S66, note~*2]{Popper2002LogicDiscoveryUI}.

Section~66 analyses the logical form with classical deducibility and
contradiction only, without a special probability logic. Probability
estimates are neither verifiable nor falsifiable. They can still have
existential consequences: an estimate with \(p\neq0,1\) entails that
ones and zeros occur, and that after every element a one and a zero
occur. These are \emph{universalized existential statements}, which are
better or worse confirmed as more or fewer of their existential
instances are verified. The requirement of randomness has this
existential form: for every \(\varepsilon>0\), every predecessor
\(n\)-tuple and every \(x\) there is a selected \(y>x\) at which the
frequency deviates from \(p\) by less than \(\varepsilon\). The
requirement of uniqueness is universal and has no extensional
consequences. It is not idle, however: it gives probability statements
the form of a universal statement conjoined with an existential
hypothesis, so that they can contradict one another. The convergence
axiom is also universal, but has only an intensional significance, the
exclusion of rule-defined sequences without limits
\parencite[\S66 and notes~2--4]{Popper2002LogicDiscoveryUI}. A later
note concedes that the probabilistic notions \emph{almost deducible} and
\emph{almost contradictory} are useful here
\parencite[\S66, note~*1]{Popper2002LogicDiscoveryUI}.

Section~67 shows the danger. Unlimited use of probability estimates can
explain any regularity, the law of gravity included, as a cosmic period
of accumulated accidents in a long enough random process. Such
explanations are unfalsifiable and are classed as speculative
metaphysics \parencite[\S67]{Popper2002LogicDiscoveryUI}. A later note
cites Haldane's argument that a probabilistic theory of entropy makes it
almost certain the world will wind itself up again
\parencite[\S67, note~*1]{Popper2002LogicDiscoveryUI}.

Section~68 confronts a physicist's definition of probability as the
value approached by relative frequencies in long series. Three objections
follow: it is circular, because only \emph{almost all} long segments
are stable; it gives no criterion of length; and it gives no criterion
of approximation. Popper keeps the definition by a methodological
decision, in his words: ``never to explain physical effects, i.e.
reproducible regularities, as accumulations of accidents''
\parencite[\S68]{Popper2002LogicDiscoveryUI}. The decision narrows the
concept of probability. For a macro effect of value \(p\) within measurement
imprecision \(\pm\varphi\), produced by \(n\) micro events, Bernoulli's
theorem relates the complementary probability \(\varepsilon\), \(n\) and
\(\Delta p\). Choosing \(\delta\) close to \(\varphi\) leaves the
question of when \(\varepsilon\) is small. The physicist's reply is to
choose \(n\) so large that \(\Delta p\) is insensitive to changes in the
order of magnitude of \(\varepsilon\). Rare deviations are not denied;
they are not reproducible effects. \emph{Reproducible} deviations from a
macro effect deduced from a probability estimate falsify the estimate.
The rule forbids the predictable and reproducible occurrence of
systematic deviations, not atypical segments as such. Popper's summary
is that probability statements are metaphysical so far as they are not
falsifiable, and are used as falsifiable statements so far as they are
used empirically \parencite[\S68]{Popper2002LogicDiscoveryUI}.
Item~3 of \cref{thm:p15-w02b-p15-027-bernoulli-bound} gives a
conservative quantitative version of the choice of \(n\). Popper's
remark that every \(n\) has a characteristic value of \(\Delta p\)
insensitive to \(\varepsilon\) is not proved here. A later note says this
analysis is clarified and superseded by appendix~*ix: large samples tend
to give a hypothesis either a high negative degree of corroboration,
which may be decided to count as falsification, or a high positive one
\parencite[\S68, note~*4]{Popper2002LogicDiscoveryUI}. Appendix~*ix
itself is not reconstructed here.

Section~69 locates chance in our inability to predict. Dice rules (the
\emph{frame conditions}) are chosen to prevent measurement of initial
conditions; that a sequence is chance-like does not show that its
elements are unpredictable, and corroborated frequency estimates do not
show that the sequence lacks law. Chance-like sequences can be
constructed by mathematical rule. Success of probability predictions is a
symptom only of the absence of \emph{simple} laws in the structure of the
sequence \parencite[\S69]{Popper2002LogicDiscoveryUI}. A later note
says this is no concession to a subjective interpretation of
probability, and another records that the metaphysical view dismissed
there is recommended in the \emph{Postscript}
\parencite[\S69, notes~*1, *2]{Popper2002LogicDiscoveryUI}.
Section~70 denies that macro laws follow from non-statistical micro laws
alone: statistical conclusions need statistical premises, such as
frequency assumptions about recurring initial conditions. Equal-chance
assumptions are not \emph{a priori}, as the differences between
classical, Bose--Einstein and Fermi--Dirac statistics show. The waterfall
illustration contrasts deducing one individual fluctuation from initial
conditions with explaining a recurring statistical distribution
\parencite[\S70]{Popper2002LogicDiscoveryUI}. A later note answers
Jordan's appeal to ergodic results: they too rest on measure-theoretic
premises with equiprobability assumptions
\parencite[\S70, note~*1]{Popper2002LogicDiscoveryUI}.

\section{Formally singular probability and the quantum chapter}
\label{sec:p15-w02b-p15-027-quantum}

\begin{hbdefinition}{def:p15-w02b-p15-027-formally-singular}{Formally singular probability (\S71)}
For an occurrence \(k\) that is an element of a reference sequence
\(\alpha\) (Popper writes \(k\mathrel{\varepsilon}\alpha\), with the
copula \(\varepsilon\) for membership), the formally singular probability
that \(k\) has the property \(\beta\) is defined by
\[
 {}_{\alpha}P_k(\beta)={}_{\alpha}F(\beta)\qquad(k\mathrel{\varepsilon}\alpha).
\]
The statement is incomplete unless the reference class is fixed
\parencite[\S71]{Popper2002LogicDiscoveryUI}.
\end{hbdefinition}

An occurrence can belong to several reference classes, and its formally
singular probabilities differ accordingly, as with a person's mortality
by age group or occupation. No general rule selects the class, though
the narrowest class with adequate statistics is often suitable. The
alleged change of an event's probability once it has occurred is a
change of reference class: learning \(k\mathrel{\varepsilon}\beta\) suggests the class
\(\beta\), where \({}_{\beta}P_k(\beta)=1\), while \({}_{\alpha}P_k(\beta)\)
is unchanged. Read as an indefinite prediction, a formally singular
statement bridges to the subjective theory, provided the objective
frequency statement remains fundamental. Popper rejects, in 1934, a
reading of it as an objective singular indeterminacy of the next throw
\parencite[\S71]{Popper2002LogicDiscoveryUI}. Later notes say these
considerations are largely superseded by proving that a formal axiom
system is satisfied by the several interpretations. They also say he no
longer objects to an event hanging in the balance, now holding the
propensity interpretation as a conjecture about the world
\parencite[\S71, notes~*2--*4]{Popper2002LogicDiscoveryUI}.

Section~72 compares this with the theory of range of Bolzano, von Kries
and Waismann. Waismann proposed measuring ratios of logical ranges by
frequencies. Popper judges this workable only once a frequency theory
exists, and drops Waismann's assumption that partially overlapping ranges
are comparable. If \(\gamma\subset\beta\), the range of
\(k\mathrel{\varepsilon}\gamma\) is at most that of
\(k\mathrel{\varepsilon}\beta\), with equality only under a law
\((x)\{[x\mathrel{\varepsilon}(\alpha.\beta)]\to(x\mathrel{\varepsilon}\gamma)\}\). When that law fails and randomness is
assumed, he asserts \({}_{\alpha}F(\gamma)<{}_{\alpha}F(\beta)\)
\parencite[\S72]{Popper2002LogicDiscoveryUI}.

\begin{hbtheorem}{thm:p15-w02b-p15-027-reference-class}{proved=annex:p15-w02b-p15-027-refs}{Reference-class change and the limit of the range comparison}
\begin{enumerate}
\item Let \(\alpha\) be finite with \(0<N(\alpha.\beta)<N(\alpha)\). Then
  \({}_{\alpha.\beta}F''(\beta)=1\) and
  \({}_{\alpha.\bar\beta}F''(\beta)=0\), while \({}_{\alpha}F''(\beta)\) is
  strictly between \(0\) and \(1\) and does not depend on which element
  of \(\alpha\) is singled out. Both an element of \(\alpha.\beta\) and
  an element of \(\alpha.\bar\beta\) exist.
\item Let \(\alpha\) be the sequence of positive integers, \(\beta\) all
  positive integers and \(\gamma=\beta\setminus\{1\}\). Then
  \(\gamma\subset\beta\) and the law
  \((x)\{[x\mathrel{\varepsilon}(\alpha.\beta)]\to(x\mathrel{\varepsilon}\gamma)\}\) fails, yet both limiting
  frequencies equal \(1\).
\end{enumerate}
\end{hbtheorem}

Item~1 is the finite form of the before-and-after point. Item~2 shows
that the strict inequality of \S72 needs more than the failure of the
universal law: a finite set of exceptions leaves limiting frequencies
equal. Popper's appeal to randomness \emph{in this respect} must carry
that extra content. The finite comparison is item~7 of
\cref{thm:p15-w02b-p15-027-finite-calculus}.

Chapter~9 applies these instruments to quantum theory. It does not
question the formalism, only the logical consequences of Born's
statistical interpretation, and it leaves aside wave fields and second
quantization. Popper announces four conclusions. (1) Heisenberg's
formulae are formally singular probability statements and so, in the
end, statistical \emph{scatter relations}. (2) Measurements more precise
than the uncertainty principle allows are compatible with the formalism
and its statistical interpretation. (3) A limit of attainable precision
is therefore an additional assumption. (4) That assumption contradicts
the statistically interpreted formulae
\parencite[chapter~9, opening]{Popper2002LogicDiscoveryUI}.
His later opening note keeps the main criticism, records a changed
interpretation of both probability and quantum theory, and excepts
\S77 as based on a mistake while singling out \S76 as important
\parencite[chapter~9, note~*1]{Popper2002LogicDiscoveryUI}.

Section~73 reports Heisenberg's programme of removing unobservable
magnitudes, such as Bohr's electron orbits, and his uncertainty
relation
\[
 \Delta x\,.\,\Delta p_x\geq\frac{h}{4\pi},
\]
read as a limit on the attainable precision of simultaneous
measurement. Popper argues that the programme was not carried out. By
combining two measurements (two positions, or a position preceded or
followed by a momentum measurement) the formalism lets one calculate a
precise path between them. Heisenberg calls the reality of that past
path a matter of taste, and Schlick calls it meaningless. Either the
particle has an exact but hidden path, a subjective reading, or it has
none although the theory calculates one, an objective reading that
leaves metaphysical elements in the formalism
\parencite[\S73]{Popper2002LogicDiscoveryUI}. A later note marks as
mistaken his 1934 claim that case (b) permits calculating the path
before the first measurement \parencite[\S73, note~4]{Popper2002LogicDiscoveryUI}.
Section~74 summarizes Born's interpretation: the squared amplitude gives
the probability of finding the particle in a region. The uncertainty
formulae are derivable from the statistically interpreted wave equation,
not the reverse, so explaining the statistical interpretation by
uncertainty reverses the relation of derivability
\parencite[\S74]{Popper2002LogicDiscoveryUI}.

Section~75 separates the \emph{Heisenberg formulae}, which follow from
the theory, from their interpretation as limits on precision, which does
not. A wave packet for one particle is a formally singular probability
statement, so the formulae derived from it are statistical. In the
\emph{statistical scatter relations}, selecting particles physically by
a position range \(\Delta x\) makes their momenta scatter over a range
\(\Delta p_x\) larger as \(\Delta x\) is smaller, and conversely. A
selection by both ranges is physically possible only when they satisfy
the displayed inequality. Popper distinguishes \emph{physical
selection} (screening off part of a beam) from measurement. Every
physical selection can serve as a measurement, but not every measurement,
for instance a Geiger counter's record of arrival, is a physical
selection. A selection made as homogeneous as the formulae permit is a
\emph{pure case}, and the scatter relations say that no aggregate is more
homogeneous than a pure case \parencite[\S75]{Popper2002LogicDiscoveryUI}.
Later notes recast the aggregate of particles as a sequence of
repetitions of an experiment with one particle or system, and restate the
pure-case thesis for experimental arrangements
\parencite[\S75, notes~*1, *2]{Popper2002LogicDiscoveryUI}.

Section~76 inverts Heisenberg's programme. Statistical formulae cannot
prohibit single events, except with probabilities 0 or 1, just as
\({}_{\alpha}P_k(\beta)=1/6\) contradicts neither
\(k\mathrel{\varepsilon}\beta\) nor \(k\mathrel{\varepsilon}\bar\beta\)
(\cref{thm:p15-w02b-p15-027-reference-class}, item~1). So alleged proofs
that exact measurement contradicts the theory fail once translated into
statistical language. What does follow is weaker: with a constructive
experimental technique, which produces initial conditions by physical
selection, the scatter relations yield frequency predictions only. The
calculable past paths are not superfluous. They are needed to test the
scatter relations, which predict that momenta scatter after a sharp
position selection; testing that needs precise single measurements of
the scattered momenta, calculated back to the slit. Popper concludes
that the statistically interpreted theory would be untestable if exact
single measurements were impossible
\parencite[\S76]{Popper2002LogicDiscoveryUI}. A later note says he still
agrees completely, and that only the precision of the prediction is
blurred, never that of a measurement. Another inserts a sentence because
he no longer accepts the \S77 argument
\parencite[\S76, notes~*1, *2]{Popper2002LogicDiscoveryUI}. A further
note withdraws his rejection of a singular objective interpretation;
Schrödinger's theory is now to be read as objective, singular and
probabilistic \parencite[\S76, note~*3]{Popper2002LogicDiscoveryUI}.

For the half-silvered mirror, the formally singular probabilities of a
photon passing or being reflected are both \(1/2\). Learning that it was
reflected seems to change them discontinuously, the reduction of the wave
packet. Popper treats this as a change of reference class from
\(\alpha\) to \(\beta\) or \(\bar\beta\), with \({}_{\alpha}P_k(\beta)\)
unchanged. Calling the logical consequence of the new information a
superluminal action is unhelpful \parencite[\S76 and note~*5]{Popper2002LogicDiscoveryUI}.
A later note says the problem had since become famous as the
discontinuous reduction of the wave packet \parencite[\S76, note~*4]{Popper2002LogicDiscoveryUI}.

Section~77 answers an opponent who grants all this but holds that
measurement and physical selection are inseparable. Popper replies that
the impossibility of exact single predictions then rests on an
additional \emph{hypothesis of linkage}, and that the quantum theory
combined with that hypothesis clashes with his view. Imagined selections
within a monochromatic beam, a slit-selected beam and a widening
packet illustrate testable path predictions for particles in small
space--time regions. He then describes crossed beams, one monochromatic
and one slit-selected, with conservation of energy and momentum. A
non-predictive measurement at \(X\) was meant to give precise predictions
for the partner particle travelling toward \(Y\), as an \emph{experimentum
crucis} against Heisenberg's interpretation
\parencite[\S77]{Popper2002LogicDiscoveryUI}.

\begin{hbwarning}{def:p15-w02b-p15-027-withdrawn-experiment}{The withdrawn experiment of \S77}
Popper's later notes state that the imaginary experiment of \S77 rests
on a mistake noted by von Weizsäcker in 1934, by Heisenberg in letters
and by Einstein in a letter reprinted in appendix~*xii, and that he has
withdrawn it. The decisive error is the claim that a momentum selection
followed by a position measurement fixes the path before the first
measurement. According to quantum theory, non-predictive measurements
fix the path only between two measurements, so one cannot know that the
particle at \(X\) came from \(P\). The last paragraph of appendix~vi
falls with it. He proposes the Einstein--Podolsky--Rosen argument as
replacement, characterizes it as weaker but valid, and states that his
arguments before the experiment are unaffected
\parencite[\S77, notes~**, *2--*4]{Popper2002LogicDiscoveryUI}.
The crossed-beam experiment is therefore not presented here as a valid
argument, and the critics' rejection of it does not by itself refute
the earlier analysis.
\end{hbwarning}

Section~78 argues that the coexistence of precision laws and frequency
hypotheses favours neither determinism nor indeterminism. Precision
statements may be derivable from frequency hypotheses as macro laws, but
frequency statements are never derivable from precision statements.
Well-corroborated frequency statements therefore do not license the
claim that there are no precise laws. Whether the world is ruled by
strict laws is metaphysical; the belief in causality hypostatizes the
methodological rule never to abandon the search for laws. No empirical
hypothesis could compel abandoning that search, since testing an
indeterminist conclusion means searching for laws. Like the bound \(c\)
on signal velocity, \(h\) forbids producing super-pure cases but does
not bar searching for them \parencite[\S78]{Popper2002LogicDiscoveryUI}.
Later notes record Einstein's opposition to the irreducibility of
frequency statements, Popper's continued agreement with the analysis, a
qualification that a falsifiable statement can have unfalsifiable
consequences, and a 1967 restatement
\parencite[\S78, notes~*1, *2, *4, *6]{Popper2002LogicDiscoveryUI}.

\section{Dated layers in the probability and quantum chapters}
\label{sec:p15-w02b-p15-027-later-notes}

The translators' note states that the 1934 text was left unchanged. New
footnotes and appendices carry starred numbers, and starred expansions of
old footnotes are marked the same way
\parencite[translators' note]{Popper2002LogicDiscoveryUI}. A starred note
therefore postdates the 1934 text. Unless it gives a year, its date is
not fixed beyond that. The table lists the starred changes to chapters~8
and~9 used above.

\begin{center}\small
\begin{tabular}{@{}p{0.17\linewidth}p{0.36\linewidth}p{0.39\linewidth}@{}}
\toprule
Locator&1934 text&Later note\\
\midrule
\S47, note~*1&Frequency theory with an axiom of randomness&Three changes: a formal calculus and a simplified frequency theory (both dated 1938; appendices~*ii--*vi), and the propensity interpretation (1953, \emph{Postscript})\\
\S53, note~4&Independence triply symmetric for finite \(\beta,\gamma\)&Condition insufficient; five-button counterexample (\cref{thm:p15-w02b-p15-027-triple-symmetry})\\
\S55, note~*2&Term \emph{free from after-effect} from Reichenbach&Added 1968: used earlier by Smoluchowski\\
\S57, note~*1&Infinite sequences with limits assumed&Finite programme not carried out; shortest \(n\)-free sequences make limits demonstrable\\
\S58, notes~*1, *2&Reformulated randomness axiom&Last seven words absent in German; Wald's construction for denumerably many systems\\
\S60, note~*3&Insensitivity to \emph{all} combinations&Restricted to those usable as gambling systems (Wald, 1935)\\
\S63, note~*1&Bernoulli's theorem to be shown without limits&Doubt about convergence kept; view implicit in the next sentence no longer held\\
\S64, note~*2&Requirement of uniqueness&May be replaced by the requirement of finitude (+2)\\
\S65, note~*1; \S66, note~*3&Probability statements not falsifiable&Shortest random-like sequences make them falsifiable; added 1972: a pointer to a solution by D.~Gillies\\
\S66, note~*2&Universalized existential form&Unbounded delay, not the form, causes unfalsifiability\\
\S68, notes~*3, *4&Reduction of other cases \emph{without much difficulty}; characteristic \(\Delta p\)&Doubted; superseded by appendix~*ix\\
\S71, notes~*2, *3&Singular objective indeterminacy rejected&Formal axioms supersede; propensity conjecture accepted\\
\S75, notes~*1, *2&Aggregates of particles&Sequences of repeated experiments\\
\S76, note~*3&Singular objective reading rejected&Objective, singular and probabilistic reading adopted\\
\S77, notes~**, *3&Crossed-beam \emph{experimentum crucis}&Based on a mistake; withdrawn\\
\bottomrule
\end{tabular}
\end{center}

The 1972 addition to note~*1 of \S65 only refers to Gillies's solution
on a later page. That solution is not reconstructed here. Nor are the new
appendices~*ii--*xii or the \emph{Postscript}, which the notes cite for
the later views.
